\section{Conclusion}
\label{sec:conclusion}

Two design choices separate NVFP4 from MXFP4, and move this paper's
backward-error proxy unequally. Across the full grid of tail weights
and contraction lengths, the scale format moves the measured backward error more
than the block size does, at every canonical cell, and the ranking survives both
secondary factors --- stochastic rounding and the random Hadamard transform ---
without reversing anywhere. The margin narrows at the heaviest tail without
crossing. Whether that ordering carries over to training loss or the
token-efficiency gap is untested. This is a post-hoc
decomposition of measured magnitudes, not a preregistered test, and it involves
no theoretical bound at any point.

That last property is what the threshold analysis leaves behind. The constructed
bound breaks across the whole tested range, so the tail index at which it first
fails is undefined everywhere on this grid, and could not carry the planned
comparison; that null motivated the pivot to comparing magnitudes directly, and
does not bear on the decomposition, which never uses the bound. Two steps would
sharpen the picture: executing the specified but unrun per-tensor reference
cells, which would place both formats against a classical baseline, and an
ablation separating the intra-block correlation mechanism from other contributors
to the $u_{\mathrm{eff}}$ residual.
